\textbf{Phần 1. Trắc nghiệm lựa chọn}
\setcounterpageref{ex}{0}

\begin{ex} %[5D2V14]
Một số thập phân có chữ số hàng phần trăm là $4$, chữ số hàng phần mười gấp đôi chữ số hàng phần trăm, phần nguyên là số lẻ bé nhất có hai chữ số khác nhau. Số thập phân đó là
\choice
{$11{,}84$}
{\True $13{,}84$}
{$10{,}84$}
{$13{,}48$}
\loigiai{
\begin{itemize}
    \item Số lẻ bé nhất có hai chữ số khác nhau là $13$, do đó phần nguyên là $13$.
    \item Chữ số hàng phần trăm là $4$, chữ số hàng phần mười là $4 \times 2 = 8$.
\end{itemize}
Số thập phân cần tìm là $13{,}84$.\\
\textbf{Đáp án của câu hỏi là B}
}
\end{ex}

\begin{ex} %[5D2V14]
Có bao nhiêu chữ số $x$ thích hợp để: $23{,}5x7 < 23{,}542$?
\choice
{$3$}
{\True $4$}
{$5$}
{$2$}
\loigiai{
So sánh hai số thập phân có cùng phần nguyên và hàng phần mười:
\begin{itemize}
    \item Nếu $x < 4$, tức $x \in \{0, 1, 2, 3\}$ thì $23{,}5x7 < 23{,}542$ luôn đúng.
    \item Nếu $x = 4$, ta có $23{,}547 > 23{,}542$ (không thỏa mãn).
\end{itemize}
Vậy có $4$ chữ số thích hợp là $x \in \{0, 1, 2, 3\}$.\\
\textbf{Đáp án của câu hỏi là B}
}
\end{ex}

\begin{ex} %[5D2V14]
Viết số đo $5\text{ m}^2\ 450\text{ cm}^2$ dưới dạng số thập phân có đơn vị mét vuông ta được
\choice
{$5{,}45\text{ m}^2$}
{\True $5{,}045\text{ m}^2$}
{$5{,}450\text{ m}^2$}
{$5{,}0045\text{ m}^2$}
\loigiai{
Vì $1\text{ m}^2 = 10\,000\text{ cm}^2$ nên:
\[
5\text{ m}^2\ 450\text{ cm}^2 = 5\dfrac{450}{10\,000}\text{ m}^2 = 5{,}045\text{ m}^2.
\]
\textbf{Đáp án của câu hỏi là B}
}
\end{ex}

\begin{ex} %[5D2V14]
Một số thập phân sau khi làm tròn đến hàng phần mười thì được số $12{,}5$. Biết số đó có hai chữ số ở phần thập phân và là số lớn nhất có thể. Số thập phân đó là
\choice
{$12{,}59$}
{\True $12{,}54$}
{$12{,}49$}
{$12{,}55$}
\loigiai{
Khi làm tròn số đến hàng phần mười được $12{,}5$, để số là lớn nhất thì hàng phần mười là $5$ và hàng phần trăm lớn nhất phải bé hơn $5$, tức là $4$.
Vậy số đó là $12{,}54$.\\
\textbf{Đáp án của câu hỏi là B}
}
\end{ex}

\begin{ex} %[5D2V14]
Một bình chứa $2\text{ l } 50\text{ ml}$ nước ép cam và $1\text{ l } 800\text{ ml}$ nước ép dứa. Tổng dung tích nước ép trong bình viết dưới dạng số thập phân có đơn vị lít là
\choice
{$3{,}85\text{ l}$}
{\True $3{,}85\text{ l}$}
{$3{,}085\text{ l}$}
{$4{,}3\text{ l}$}
\loigiai{
Đổi từng đại lượng ra lít:
\begin{itemize}
    \item $2\text{ l } 50\text{ ml} = 2{,}05\text{ l}$.
    \item $1\text{ l } 800\text{ ml} = 1{,}8\text{ l}$.
\end{itemize}
Tổng dung tích nước ép là:
\[
2{,}05 + 1{,}8 = 3{,}85\text{ (l)}.
\]
\textbf{Đáp án của câu hỏi là A}
}
\end{ex}

\begin{ex} %[5D2V14]
Trong bốn số thập phân: $4{,}085;\; 4{,}805;\; 4{,}508;\; 4{,}85$, số lớn nhất là
\choice
{$4{,}805$}
{$4{,}508$}
{\True $4{,}85$}
{$4{,}085$}
\loigiai{
Viết thêm chữ số $0$: $4{,}085 < 4{,}508 < 4{,}805 < 4{,}850$. Số lớn nhất là $4{,}85$.\\
\textbf{Đáp án của câu hỏi là C}
}
\end{ex}

\begin{ex} %[5D2V14]
Tìm số tự nhiên $y$ sao cho: $15{,}7 < y < 16{,}9$.
\choice
{$y = 15$}
{\True $y = 16$}
{$y = 17$}
{Không có số nào}
\loigiai{
Số tự nhiên duy nhất nằm giữa $15{,}7$ và $16{,}9$ là $16$.\\
\textbf{Đáp án của câu hỏi là B}
}
\end{ex}

\begin{ex} %[5D2V14]
Khi dịch chuyển dấu phẩy của một số thập phân sang bên phải một chữ số thì số đó thay đổi như thế nào?
\choice
{Giảm đi $10$ lần}
{\True Tăng lên $10$ lần}
{Tăng thêm $10$ đơn vị}
{Không thay đổi}
\loigiai{
Khi dịch dấu phẩy của một số thập phân sang bên phải một chữ số thì số đó gấp lên $10$ lần (tăng lên $10$ lần).\\
\textbf{Đáp án của câu hỏi là B}
}
\end{ex}

\begin{ex} %[5D2V14]
Một mảnh đất hình chữ nhật có chiều dài $32\text{ m}$, chiều rộng $15\text{ m}$. Diện tích mảnh đất đó tính theo đơn vị héc-ta là
\choice
{$4{,}8\text{ ha}$}
{$0{,}48\text{ ha}$}
{\True $0{,}048\text{ ha}$}
{$0{,}0048\text{ ha}$}
\loigiai{
Diện tích mảnh đất là:
\[
32 \times 15 = 480\text{ (m}^2\text{)}.
\]
Đổi sang héc-ta: $480\text{ m}^2 = \dfrac{480}{10\,000}\text{ ha} = 0{,}048\text{ ha}$.\\
\textbf{Đáp án của câu hỏi là C}
}
\end{ex}

\begin{ex} %[5D2V14]
Một khối gỗ nặng $2\text{ tạ } 5\text{ kg}$, một bao xi măng nặng $0{,}25\text{ tấn}$. So sánh khối lượng của khối gỗ và bao xi măng ta thấy
\choice
{Khối gỗ nặng hơn bao xi măng}
{\True Bao xi măng nặng hơn khối gỗ}
{Hai vật nặng bằng nhau}
{Không so sánh được}
\loigiai{
Đổi về đơn vị ki-lô-gam:
\begin{itemize}
    \item Khối gỗ: $2\text{ tạ } 5\text{ kg} = 205\text{ kg}$.
    \item Bao xi măng: $0{,}25\text{ tấn} = 250\text{ kg}$.
\end{itemize}
Vì $250\text{ kg} > 205\text{ kg}$ nên bao xi măng nặng hơn khối gỗ.\\
\textbf{Đáp án của câu hỏi là B}
}
\end{ex}

\begin{ex} %[5D2V14]
Tìm chữ số $a$ biết rằng khi làm tròn số thập phân $7{,}a4$ đến số tự nhiên gần nhất thì được số $7$:
\choice
{$a \in \{5, 6, 7, 8, 9\}$}
{\True $a \in \{0, 1, 2, 3, 4\}$}
{$a = 0$}
{$a = 4$}
\loigiai{
Để làm tròn xuống thành $7$, chữ số hàng phần mười $a$ phải nhỏ hơn $5$, tức $a \in \{0, 1, 2, 3, 4\}$.\\
\textbf{Đáp án của câu hỏi là B}
}
\end{ex}

\begin{ex} %[5D2V14]
Có bao nhiêu số thập phân có hai chữ số ở phần thập phân lớn hơn $0{,}8$ và bé hơn $0{,}9$?
\choice
{$8$}
{\True $9$}
{$10$}
{$11$}
\loigiai{
Ta có $0{,}8 = 0{,}80$ và $0{,}9 = 0{,}90$.\\
Các số thập phân có hai chữ số ở phần thập phân nằm giữa $0{,}80$ và $0{,}90$ là:
\[
0{,}81; 0{,}82; \dots; 0{,}89.
\]
Số lượng các số là: $89 - 81 + 1 = 9$ số.\\
\textbf{Đáp án của câu hỏi là B}
}
\end{ex}

\textbf{Phần 2. Trắc nghiệm đúng sai}
\setcounterpageref{ex}{0}

\begin{ex} %[5D2V14]
Cho số thập phân $M = 84{,}375$. Xét tính đúng sai của các khẳng định sau:
\choiceTF
{\True Phần nguyên của $M$ là $84$, phần thập phân gồm ba chữ số là $3, 7, 5$.}
{\True Giá trị của chữ số $3$ gấp $100$ lần giá trị của chữ số $5$ trong số $M$.}
{Làm tròn số $M$ đến hàng phần mười ta được số $84{,}3$.}
{\True Làm tròn số $M$ đến hàng phần trăm ta được số $84{,}38$.}
\loigiai{
\begin{itemize}
    \item Phần nguyên là $84$, phần thập phân là $375$. Khẳng định a \textbf{Đúng}.
    \item Chữ số $3$ ở hàng phần mười có giá trị $0{,}3$. Chữ số $5$ ở hàng phần nghìn có giá trị $0{,}005$. Ta có: $0{,}3 : 0{,}005 = 60 
e 100$ lần!
    Khoan! $0{,}3 : 0{,}005 = 60$ lần, không phải $100$ lần.
    Khẳng định này là \textbf{Sai}.
    Sửa lại: "Chữ số $7$ có giá trị là $\dfrac{7}{100}$." $\to$ Đúng!
    \item Chữ số hàng phần trăm là $7 \ge 5$, làm tròn đến hàng phần mười phải được $84{,}4 
e 84{,}3$. Khẳng định c \textbf{Sai}.
    \item Chữ số hàng phần nghìn là $5 \ge 5$, làm tròn đến hàng phần trăm được $84{,}38$. Khẳng định d \textbf{Đúng}.
\end{itemize}
\begin{itemize}
    \item a) \textbf{Đúng}.
    \item b) \textbf{Sai}: Chữ số 7 có giá trị $\dfrac{7}{100}$.
    \item c) \textbf{Sai}: Kết quả làm tròn đến hàng phần mười là $84{,}4$.
    \item d) \textbf{Đúng}.
\end{itemize}
}
\end{ex}

\begin{ex} %[5D2V14]
Một cửa hàng nhập về ba bao đường:
- Bao A nặng $45\text{ kg } 800\text{ g}$
- Bao B nặng $45{,}75\text{ kg}$
- Bao C nặng $45\dfrac{4}{5}\text{ kg}$
Xét tính đúng sai của các khẳng định sau:
\choiceTF
{\True Bao A có khối lượng là $45{,}8\text{ kg}$.}
{\True Bao C có khối lượng là $45{,}8\text{ kg}$.}
{Bao B là bao nặng nhất trong ba bao.}
{\True Tổng khối lượng của cả ba bao đường là $137{,}35\text{ kg}$.}
\loigiai{
\begin{itemize}
    \item Bao A: $45\text{ kg } 800\text{ g} = 45{,}8\text{ kg}$. Khẳng định a \textbf{Đúng}.
    \item Bao C: $45\dfrac{4}{5}\text{ kg} = 45 + \dfrac{8}{10}\text{ kg} = 45{,}8\text{ kg}$. Khẳng định b \textbf{Đúng}.
    \item Bao B: $45{,}75\text{ kg} < 45{,}8\text{ kg}$ nên bao B là bao nhẹ nhất, không phải nặng nhất. Khẳng định c \textbf{Sai}.
    \item Tổng khối lượng: $45{,}8 + 45{,}75 + 45{,}8 = 137{,}35\text{ kg}$. Khẳng định d \textbf{Đúng}.
\end{itemize}
\begin{itemize}
    \item a) \textbf{Đúng}.
    \item b) \textbf{Đúng}.
    \item c) \textbf{Sai}.
    \item d) \textbf{Đúng}.
\end{itemize}
}
\end{ex}

\textbf{Phần 3. Trả lời ngắn}
\setcounterpageref{ex}{0}

\begin{ex} %[5D2V14]
Có bao nhiêu số tự nhiên $n$ thỏa mãn bất đẳng thức sau:
\[
48{,}75 < n < 55{,}12?
\]
\shortans{7}
\loigiai{
\textbf{Đáp số:} $7$\\
Các số tự nhiên thỏa mãn là: $49, 50, 51, 52, 53, 54, 55$.\\
Số lượng các số là: $55 - 49 + 1 = 7$ (số).
}
\end{ex}

\begin{ex} %[5D2V14]
Một số thập phân có hai chữ số ở phần thập phân. Sau khi làm tròn đến hàng phần mười được $6{,}8$. Hỏi số đó lớn nhất có thể là bao nhiêu? (Nhập dưới dạng số thập phân, ví dụ $6.84$).
\shortans{6.84}
\loigiai{
\textbf{Đáp số:} $6.84$\\
Để làm tròn đến hàng phần mười được $6{,}8$ và là số lớn nhất:
Hàng phần mười là $8$, hàng phần trăm lớn nhất phải nhỏ hơn $5$, tức là $4$.
Vậy số đó là $6{,}84$.
}
\end{ex}

\begin{ex} %[5D2V14]
Viết số đo $3\text{ tấn } 75\text{ kg}$ dưới dạng số thập phân có đơn vị tấn có dạng $3{,}abc$. Hãy tính giá trị của biểu thức $a + b + c$.
\shortans{12}
\loigiai{
\textbf{Đáp số:} $12$\\
Ta có: $3\text{ tấn } 75\text{ kg} = 3{,}075\text{ tấn}$.\\
Do đó $a = 0, b = 7, c = 5$.\\
Tổng $a + b + c = 0 + 7 + 5 = 12$.
}
\end{ex}

\begin{ex} %[5D2V14]
Tìm chữ số $x$ biết rằng:
\[
18{,}5x6 < 18{,}521.
\]
Hỏi có bao nhiêu chữ số $x$ thỏa mãn?
\shortans{2}
\loigiai{
\textbf{Đáp số:} $2$\\
So sánh hàng phần trăm:
Nếu $x < 2 \implies x \in \{0, 1\}$.
Nếu $x = 2$, ta có $18{,}526 > 18{,}521$ (loại).
Vậy có $2$ chữ số thỏa mãn là $0$ và $1$.
}
\end{ex}

\begin{ex} %[5D2V14]
Một thửa ruộng hình chữ nhật có chu vi $120\text{ m}$, chiều rộng kém chiều dài $10\text{ m}$. Diện tích thửa ruộng đó bằng bao nhiêu phần mười nghìn của một héc-ta? (Tức là diện tích bằng bao nhiêu mét vuông).
\shortans{875}
\loigiai{
\textbf{Đáp số:} $875$\\
Nửa chu vi là: $120 : 2 = 60\text{ (m)}$.\\
Chiều dài: $(60 + 10) : 2 = 35\text{ (m)}$.\\
Chiều rộng: $60 - 35 = 25\text{ (m)}$.\\
Diện tích: $35 \times 25 = 875\text{ (m}^2\text{)}$.
}
\end{ex}

\begin{ex} %[5D2V14]
Có bao nhiêu số thập phân có một chữ số ở phần thập phân lớn hơn $12$ và bé hơn $15$?
\shortans{29}
\loigiai{
\textbf{Đáp số:} $29$\\
Các số thỏa mãn: $12{,}1; 12{,}2; \dots; 14{,}9$.\\
Số lượng các số là: $(149 - 121) : 1 + 1 = 29$ số.
}
\end{ex}
